%%=======================================
%  WHUT Bachelor Thesis v 0.1b
%  
%  Tsao Yu
%  tsaoyu@tsaoyu.com
%  Last Modify Date Feb.6,2014
%%=======================================
\documentclass[cs4size,a4paper]{ctexart}
%==================== 数学符号公式 ============
\usepackage{xeCJK}
\newcommand{\zhongsong}{\CJKfontspec{STZhongsong}}%华文中宋，请自行下载字体并安装

\usepackage{amsmath}                 % AMS LaTeX宏包
%\usepackage{amssymb}                % 用来排版漂亮的数学公式
%\usepackage{amsbsy}
\usepackage[style=1]{mdframed}
\usepackage{amsthm}
\usepackage{amsfonts}
\usepackage{mathrsfs}                % 英文花体字 体
\usepackage{bm}                      % 数学公式中的黑斜体
\usepackage{bbding,manfnt}           % 一些图标，如 \dbend
\usepackage{lettrine}                % 首字下沉，命令\lettrine
\usepackage[backend=biber,style=gb7714-2015]{biblatex}  %配置gb7714-2015引用格式，使用biber编译，自行安装biber
\addbibresource{bibfile.bib}                            %添加bib file
\def\attention{\lettrine[lines=2,lraise=0,nindent=0em]{\large\textdbend\hspace{1mm}}{}}
\usepackage{longtable}
\usepackage[toc,page]{appendix}
\usepackage{geometry}                % 页边距调整
\geometry{top=2.5cm,bottom=2cm,left=2.5cm,right=2cm}
%\usepackage{relsize}                % 调整公式字体大小：\mathsmaller,\mathlarger
%\usepackage{caption2}               % 浮动图形和表格标题样式
\usepackage{booktabs}                %三线表上下加粗
\usepackage{diagbox}                 % 分类表头
%%%使用带圈数字作为教主
\usepackage{pifont}
\usepackage[symbol*,stable]{footmisc}
\DefineFNsymbols{circled}{{\ding{192}}{\ding{193}}{\ding{194}}
{\ding{195}}{\ding{196}}{\ding{197}}{\ding{198}}{\ding{199}}{\ding{200}}{\ding{201}}}
\setfnsymbol{circled}
%====================公式按章编号==========================
\numberwithin{equation}{section}
\numberwithin{table}{section}
\numberwithin{figure}{section}
%================= 基本格式预置 ===========================
\usepackage{fancyhdr}
\pagestyle{fancy}
\fancyhf{}  
\fancyhead[C]{\zihao{5}  \songti 武汉理工大学毕业设计（论文）}
\fancyfoot[C]{~\zihao{5} \thepage~}
\renewcommand{\headrulewidth}{0.65pt} 
\CTEXsetup[format={\centering\bfseries\zihao{-2}},name={第, 章}]{section}
\CTEXsetup[nameformat={\bfseries\zihao{3}}]{subsection}
\CTEXsetup[nameformat={\bfseries\zihao{4}}]{subsubsection}
%================== 图形支持宏包 =========================
\usepackage{subfigure}
\usepackage{graphicx}                % 嵌入png图像
\usepackage{color,xcolor}            % 支持彩色文本、底色、文本框等
\usepackage{hyperref}                % 交叉引用
\usepackage{caption}

% set up labelformat and labelsep for figure
\captionsetup{labelsep=quad}
\captionsetup{figurewithin=section}

\renewcommand{\thesubfigure}{(\arabic{subfigure})} %还可设置图编号显示格式，加括号或者不加括号

%==================== 源码和流程图 =====================
\usepackage{listings}                % 粘贴源代码
\usepackage{tikz}                    
\usepackage{tikz-3dplot}
\usetikzlibrary{shapes,arrows,positioning}
%===================   正文开始    ===================
\begin{document}
%\bibliographystyle{gbt7714-2005}     %论文引用格式
%===================  定理类环境定义 ===================
\newtheorem{example}{例}              % 整体编号
\newtheorem{algorithm}{算法}
\newtheorem{theorem}{定理}            % 按 section 编号
\newtheorem{definition}{定义}
\newtheorem{axiom}{公理}
\newtheorem{property}{性质}
\newtheorem{proposition}{命题}
\newtheorem{lemma}{引理}
\newtheorem{corollary}{推论}
\newtheorem{remark}{注解}
\newtheorem{condition}{条件}
\newtheorem{conclusion}{结论}
\newtheorem{assumption}{假设}
%==================重定义 ===================
\renewcommand{\contentsname}{目~~录}     
\renewcommand{\abstractname}{摘~~要} 
\renewcommand{\refname}{参考文献}     
\renewcommand{\indexname}{索引}
\renewcommand{\figurename}{图}
\renewcommand{\tablename}{表}
\renewcommand{\appendixname}{附录}
\renewcommand{\proofname}{证明}
\renewcommand{\algorithm}{算法} 
%============== 封皮和前言 =================
%===============  封面  =================
\smallskip
\begin{center}
\begin{figure}[!th]
\centering
\includegraphics[width=0.7\linewidth]{figure/SchoolName}
\end{figure}

\vspace*{1.0cm}
{\zihao{1} 毕业设计（论文）} \\
\vspace*{2.1cm}
 \heiti{\zihao{1} 武汉理工本科论文\LaTeX 模板 }\\
\vspace*{5.0cm}
\songti
\begin{tabular}{cc}
 \zihao{-2} 学院（系）：&\underline{\makebox[7cm][c]{\zihao{-2}交通学院}} \\ 
 \\
 \zihao{-2}专业班级: & \underline{\makebox[7cm][c]{\zihao{-2}船舶与海洋工程1006班}} \\ 
 \\
 \zihao{-2}学生姓名: & \underline{\makebox[7cm][c]{\zihao{-2}曹宇}} \\ 
 \\
 \zihao{-2}指导教师: & \underline{\makebox[7cm][c]{\zihao{-2}徐海祥}} \\ 
 \\
\end{tabular} 
\end{center}
\thispagestyle{empty}
\clearpage
%=====================原创性声明===========
\begin{center}
\zihao{-2} \textbf{\songti 学位论文原创性声明}
\end{center}

本人郑重声明：所呈交的论文是本人在导师的指导下独立进行研究所取得的研究成果。除了文中特别加以标注引用的内容外，本论文不包括任何其他个人或集体已经发表或撰写的成果作品。本人完全意识到本声明的法律后果由本人承担。 
\begin{flushright}
\zihao{4} 作者签名：\qquad ~~~\\

年\qquad 月\qquad 日
\end{flushright}
\vskip 2cm
\begin{center}
\zihao{-2} \textbf{\songti 学位论文版权使用授权书}
\end{center}

本学位论文作者完全了解学校有关保障、使用学位论文的规定，同意学校保留并向有关学位论文管理部门或机构送交论文的复印件和电子版，允许论文被查阅和借阅。本人授权省级优秀学士论文评选机构将本学位论文的全部或部分内容编入有关数据进行检索，可以采用影印、缩印或扫描等复制手段保存和汇编本学位论文。\smallskip

本学位论文属于
\begin{tabular}[t]{l}
1、保密$ \Box$，在~~~年解密后适用本授权书  \\ 
2、不保密$ \Box$  \\ 
\end{tabular} \\
\begin{center}
（请在以上相应方框内打“$\surd”$）
\end{center}
\begin{flushright}
\zihao{4} 作者签名：  \quad\quad\quad\quad 年 \quad  月  \quad  日\\
导师签名：   \quad\quad\quad\quad 年 \quad  月 \quad   日\\
\end{flushright}
\thispagestyle{empty}
\clearpage

%=============设计（论文）任务书===========
\begin{center}
\zihao{-2}\textbf{\songti 本科生毕业设计（论文）任务书} 
\end{center}
\smallskip
\renewcommand{\arraystretch}{1.3}
\begin{tabular}{lll}
\zihao{4} \textbf{\songti 学生姓名: 曹宇} & & \zihao{4} \textbf{\songti 专业班级:\quad\quad 船海1006班} \\ 
\zihao{4} \textbf{\songti 指导教师：徐海祥}&\makebox [3cm] & \zihao{4} \textbf{\songti 工作单位：\quad 武汉理工大学} \\ 
\end{tabular}\\
\begin{tabular}{lll}
\zihao{4} \textbf{\songti 设计(论文)题目:}& \zihao{4} \textbf{\songti  武汉理工本科论文\LaTeX 模板 } &\\ 
\zihao{4} \textbf{\songti 设计（论文）主要内容：} \\
\end{tabular} \\ 
\begin{enumerate}
\item \LaTeX 环境的配置
\item 主要字体的控制和数学公式的选用
\item 图表和代码的粘贴
\end{enumerate}
\begin{tabular}{ll}
\zihao{4} \textbf{\songti 要求完成的主要任务:}
\end{tabular} \\ 
\begin{enumerate}
\item 选择合适的\TeX 编辑系统
\item 学习如何使用控制代码完成排版
\item 合理的安排学习和科研的时间来发展自己兴趣爱好
\end{enumerate}
\begin{tabular}{ll}
\zihao{4} \textbf{\songti 必读参考资料：}
\end{tabular}
\begin{enumerate}
\item \LaTeX  \quad User Manual
\item  字体设计的艺术
\end{enumerate}
\begin{tabular}{lll}
\zihao{4} \textbf{\songti 指导教师签名： }&\makebox [4cm]& \zihao{4} \textbf{\songti 系主任签名：} \\
& & \zihao{4} \textbf{\songti 院长签名（章）}
\end{tabular}
\thispagestyle{empty}
\clearpage
%==========本科生毕业设计（论文）开题报告  =============
\begin{center}
\zihao{-2} \textbf{\songti 武汉理工大学}\\
\zihao{-2} \textbf{\songti 本科生毕业设计（论文）开题报告} 
\end{center}
\begin{tabular}{|l|}
\hline \rule[-2ex]{0pt}{5.5ex} \makebox[13.5cm][l]{\zihao{4} \heiti 1、目的及意义（含国内外的研究现状分析） } \\ 
\quad \LaTeX 是国际通行的科技论文排版软件，国际上科研机构和大学都采用它写作\\
\quad 国内著名高校都有自己的本科生\LaTeX 模板供毕业生使用\\
\quad 但是武汉理工大学还没有本科生\LaTeX 模板可以参考\\
\quad 人类的价值在于创造而不是索取 \\
\hline \rule[-2ex]{0pt}{5.5ex}  \zihao{4} \heiti
2、基本内容和技术方案\\ 
\quad 采用GITHUB托管降低代码维护成本\\
\quad 加入在线\TeX 编辑器的使用简介 \\
\quad 授人以渔，注重方法和理念的引导\\
\hline \rule[-2ex]{0pt}{5.5ex}  \zihao{4} \heiti
3、进度安排 \\ 
\quad 离 deadline 两个月吃喝玩乐 \\
\quad 离 deadline 一个月吃喝玩乐 \\
\quad 离 deadline 半个月吃喝玩乐 \\
\quad 离 deadline 一个星期狂写论文 \\
\hline \rule[-2ex]{0pt}{5.5ex} \zihao{4} \heiti
4、指导教师意见 \\ 
\quad 曹宇同学是个好同志\\
\quad 曹宇同志是个好同学\\
\quad 本表格是支持跨页的长表格，你可以复制上面的内容进行测试\\
\quad 具体方法是将tabular改为 longtable然后再编译\\
\makebox[10cm][r]指导教师签名：\\
\makebox[12cm][r]\quad 年\quad 月\quad 日\\
\hline 
\end{tabular} 
\thispagestyle{empty}
\clearpage
\pagestyle{plain}
\pagenumbering{Roman}
\section*{\zihao{2} \centering 摘要}

\vskip0.5cm
本文基于武汉理工大学本科生毕业论文格式2013版的相关要求，结合\LaTeX 在实际运用中的基本技巧和方法对于科技论文排版方法进行一个简略的介绍。通过参照本科生毕业论文的相关要求，实现了符合国际科技论文排版规则的具有一定美感的毕业论文模板设计。 


\textbf{关键词:}  动力定位， 船舶操纵性 ，控制方法，状态估计算法
\addcontentsline{toc}{section}{摘要}

\clearpage
\section*{\zihao{2} \centering \textbf{Abstract} }
   %用了Times New Roman字体来美化观感

In this short article we will discuss about \LaTeX\,  for your dissertation \\

\textbf{Key Words:} Dynamic Positioning, Ship Manoeuvrability ,Control Algorithm, State Estimate Algorithm
\addcontentsline{toc}{section}{Abstract}





\pagestyle{empty}
\tableofcontents 
\thispagestyle{empty}
%============== 论文正文   =================
\pagestyle{fancy}
\section{基本编译环境}
\subsection{编辑环境}

本模板采用\TeX Studio 2.6.6 作为编辑器，\TeX Live 2013套件作为编译器。
他们都是跨平台的编译软件，更重要的是他们都是完全免费的软件。 安装的顺序是先安装
\TeX Live 2013套件，再安装\TeX Studio 编辑器。 \TeX Studio编辑器的优点
有：
\begin{enumerate}
\item 清晰的组织结构，你可以在屏幕左侧看到他们
\item 便捷的自动补充功能，只要输入命令的一部分就能够完成撰写
\item 合理的宏包查看方式，右键菜单中可以找到宏包的文档
\item 贴心的实用工具，矩阵插入助手，表格编辑助手等
\end{enumerate}

当然还有其他的编译方式，比如WinEdt,LyX,Vim，甚至在线的编辑器。我也会把这份样板
上传到ShareLaTeX上，这样你就可以方便的使用在线编辑系统了。但不论你使用哪一种编辑软件或者编译器，请务必拥有他们的说明手册。这里，\LaTeX 是使用
一个个各自独立的宏作为功能的概念的，所以这些说明手册就是宏包的文档。
比如本模板使用了美国数学协会的\AmS math宏包来完成数学公式的编写，在你需要
编写数学公式你就需要找到\AmS math宏包的说明文档，文档中会告诉你这个宏包如何
使用控制语句编写数学公式。这些文档可以在网络上，或者下载的套件中找到。当然，最便捷的途径已经在上面的features里面告诉你了。
\subsection{如何使用本模板}
当你获得了所有必需的软件之后，就可以试着编译这个文档了。编译的要求是把thesis.tex设定为主文档，默认编辑器（Default compiler）则使用Xe\LaTeX 。在\TeX Studio 中的设置方法是点击Option之后选择Setup，把Build选项卡里的Default compiler 改为Xe\LaTeX 即可。如果你在编译过程中遇到红色字体提示的错误，
请拷贝错误的内容到Google中寻找解决他们的方法。此时需要你有着一定的英文阅读
能力，因为大部分\LaTeX 的文档还仍是英文的。

如果你还需要有关\LaTeX 的快速入门书籍，首先推荐阅读宏包的帮助文档，其次可以选择一些短小（200页以内）的入门小册子，中文书籍首推北京大学刘海洋编写的《\LaTeX 入门》。相信这些文档和资料可以帮助你快速的入门，发现自己所欠缺的内容之后逐渐学习最终成为\LaTeX 排版的专家。       %
\section{基本文档功能的使用}
\subsection{字体字号的控制}
对于中国用户来说， 学习\LaTeX 最大的困难首先不是阅读各种各样的宏包文档
而是需要面对\LaTeX 中一个最为困难的问题：实现汉字的显示和控制。在
\LaTeX 创始的时候Kunth老爷子没有想到这世界上还有这么复杂的文字，好在我们
的时代进步。加之无数前人的努力，开发出了一套适合中文汉字显示的\LaTeX 宏包，本论文模板正是使用这个宏包作为开发的基础。

普通大概会花上两三天的时间来纠结中文显示的问题，这一部分是\LaTeX 本身的不完美造成的，另一部分是由于我们盲目的查询，复制，粘贴网络上来源不明的代码造成的。这样做非常容易造成宏包间的冲突，从而导致编译失败。学习\LaTeX 的最好方法是从最简单的样式开始，逐渐接触模板的控制方法这样才能少走弯路。
\subsubsection{字体}
根据中文汉字支持宏包ctexart的参考文档，模板中预置的常用字体一共有五种，他们分别是：宋体，黑体，仿宋，楷书， 隶书。对应的控制方式如下：
\begin{center}
\begin{tabular}{ccccc}
\hline \rule[-2ex]{0pt}{5.5ex} { 宋体} & { 黑体} & { 仿宋} & { 隶书} & {楷书} \\ 
\hline \rule[-2ex]{0pt}{5.5ex} \textbackslash songti &\textbackslash  heiti  & \textbackslash fangsong & \textbackslash lishu & \textbackslash kaishu \\ 
\hline 
\end{tabular}  
\end{center}
这些字体基本满足了武汉理工大学本科生毕业论文中所要求的字体的需求，当然如果你有更多的字体需求可以参照ctexart的参考文档进行进一步的设置。Linux用户请用
Adobe的免费字体进行替换，当然对于Linux用户来说这些都不算问题,Macintosh 用户请用类似的苹果系统字体进行替换。
\subsubsection{字号}
大部分的字号不需要用户自行调整，所有的标题格式和正文格式已经完成了预置。为了防止特殊情况，在此对于字号的控制方法进行简单的介绍。\LaTeX 中的字体定义是pt作为单位的，而Word中的要求则大部分以小一号这样的方式定义的。通常的，小四号对应12pt，亦是本模板正文中所使用的字号。中文字号的设置命令是\textbackslash zihao\{字号\}, 在字号前加-（负号）则代表小一号。如 \textbackslash zihao\{4\} 代表四号字， \textbackslash zihao\{-4\} 则代表小四号。建议用户在正文不要使用这些命令，更不要使用这些命令来完成格式的调整。
\subsection{数学公式和证明}
\subsubsection{ 插入数学公式}
\LaTeX 之父 Stanford教授 Knuth 用\$ 符号界定数学公式，暗指着每个好的公式都是无价之宝。正是\TeX 系统让这些公式变得美丽而和谐，同时输入数学公式变成简单愉快的过程。
不过你需要的基本是\AmS math宏包交给你的，列举我最喜欢的一个公式。
如，
\[  e^{i\pi}+1=0\]
下面举例如何使用预置的文献撰写格式，此部分原作者是北京大学计算机系的于江生。
\begin{theorem}[L\'{e}vy\index{L\'{e}vy 定理}]
令 $F(x),\varphi(t)$ 分别为随机变量 $X$ 的分布函数和特征函数。
假定 $F(x)$ 在 $a+h$ 和 $a-h (h>0)$ 处连续，则有
\begin{align}
 \label{Levy theorem}  % 方程的标记可以是专有名词
F(a+h)-F(a-h)&=\lim_{T\rightarrow\infty}\frac{1}{\pi}\int^{T}_{-T}\frac{\sin ht}{t} 
e^{-ita}\varphi(t)dt
\end{align}
\end{theorem}
\begin{proof}
  从略。感兴趣的读者可以参考……。
\end{proof}


\begin{corollary}
密度函数和特征函数之间有如下的关系。
\begin{align}
 \label{DensityCharacteristic}   % 自定义的标记
  f(x)&=\frac{1}{2\pi}\int^{+\infty}_{-\infty} e^{-itx}\varphi(t)dt
\end{align}
\end{corollary}

\begin{proof}
由公式 (\ref{Levy theorem}) 和 Lebesgue 定理，我们有
\begin{align*}
 \frac{F(x+\Delta x)-F(x)}{\Delta x}&=\frac{1}{2\pi}\int^{+\infty}_{-\infty}
 \frac{\sin(t\Delta x/2)}{t\Delta x/2} e^{-it(x+\Delta x/2)}\varphi(t) dt\\
  f(x)&=\frac{1}{2\pi}\int^{+\infty}_{-\infty}\lim_{\Delta x\rightarrow 0}
 \frac{\sin(t\Delta x/2)}{t\Delta x/2} e^{-it(x+\Delta x/2)}\varphi(t) dt\\
  &=\frac{1}{2\pi}\int^{+\infty}_{-\infty} e^{-itx}\varphi(t)dt\qedhere
\end{align*}
\end{proof}

我们知道特征函数的定义是

\begin{align}
 \label{section1:characteristic}   % 标记中注明了章节号
 \varphi(t)&= E(e^{itX})=\int^{+\infty}_{-\infty} e^{itx} f(x)dx
\end{align}

L\'{e}vy 定理在分布函数和特征函数之间搭建了一座桥梁。
对比 (\ref{DensityCharacteristic}) 和 (\ref{section1:characteristic}) 可见，
密度函数和特征函数之间的关系非常巧妙。

                                                                                                                 
\attention 在\TeX 环境里，数学公式的表达是很自然的，绝大多数命令就是英文的数学专有名词或它们的缩写，
如果你以前读过英文的数学文献，记忆这些命令是不难的。如果你没读过，正好通过记忆这些命令来了解术语。


手头有个命令快速寻查表是很方便的，我用的是 Hypertext Help with \LaTeX ，网上可以搜到，是免费的。
\subsubsection{预置环境}
相信你已经发现了，模板中预置了有关的环境信息和这些环境下的字体样式。你只要在合适的位置，使用这些环境就能够让你文章变得简洁而美观。环境的样式如下：
\begin{center}

\begin{tabular}{c|c|c|c}
\hline\hline
环境 & 含义 & 环境& 含义\\\hline\hline

%thankspage&致谢页&code&代码\\\hline
theorem &定理&lemma &引理  \\\hline
example &例&algorithm &算法  \\\hline
definition &定义  &axiom &公理  \\\hline
property &性质  &proposition &命题 \\\hline
corollary& 推论 &remark &注解  \\\hline
condition &条件  &conclusion &结论  \\\hline
assumption &假设  &prove &证明 \\\hline
proof&证明 &&\\ \hline\hline
\end{tabular}

\end{center}
\subsection{图形表格等浮动对象}
\index{贝叶斯方法}贝叶斯方法主要用于小样本数据分析，它利用参数先验分布和
后验分布之差异进行统计推断，其一般步骤是：

\begin{enumerate}
 \item 构建概率模型，包括参数的先验分布。
 \item 给定观察数据，计算参数的后验分布。
 \item 分析模型的效果，如有必要，回到第一步。
\end{enumerate}

\begin{example}
下面，我们给一个表格的例子，一个图形的例子。

\begin{center}
\begin{table}[!ht]     % 强制在原位显示表格
\centering

\begin{tabular}{l|ccccc|c}
  $_X$\hspace{3mm} $^Y$&$y_1$&$y_2$&$\cdots$&$y_j$&$\cdots$\\
\hline
$x_1$   &$p_{11}$&$p_{12}$&$\cdots$&$p_{1j}$&$\cdots$&$p_{1\cdot}$\\
$x_2$   &$p_{21}$&$p_{22}$&$\cdots$&$p_{2j}$&$\cdots$&$p_{2\cdot}$\\
$\vdots$&$\vdots$&$\vdots$&$\vdots$&$\vdots$&$\vdots$&$\vdots$\\
$x_i$   &$p_{i1}$&$p_{i2}$&$\cdots$&$p_{ij}$&$\cdots$&$p_{i\cdot}$\\
$\vdots$&$\vdots$&$\vdots$&$\vdots$&$\vdots$&$\vdots$&$\vdots$\\
\hline
   &$p_{\cdot 1}$&$p_{\cdot 2}$&$\cdots$&$p_{\cdot j}$&$\cdots$&1

\end{tabular}
\end{table}
\end{center}

%在表\ref{marginal distribution} 中，$p_{\cdot %j}=\sum\limits_i p_{ij}$，
%类似地，$ p_{i\cdot}=\sum\limits_j p_{ij}$。
\end{example}
\subsection{插入图片}
要说\LaTeX 的缺点，其实也不是没有。图文混排能力比较差就是其中一个，相比于其他的商业办公文书软件，\LaTeX 在编辑大量的文字图形内容的时候就显得捉襟见肘了。这时用户就需要其他的专业软件了，比如Adobe公司的InDesign等专业的图文排版软件，来协助我们完成这样的任务了。这同样也是软件和用户双向选择的问题，\LaTeX 的使用人群是以数学物理作为基础的理工科人群。他的强项在于科技文献而非建筑学生的作品集，商科学生的案例答辩，希望用户了解这个常识。
通过\TeX Studio的图片助手功能，我们可以方便的完成简单的图片内容插入\footnote[1]{图片的大小，位置调整请查询graphicx宏包。}， 支持png，pdf，jpg，eps等格式的图片内容。\\
\begin{figure}[htbp] 
\centering
\includegraphics {figure/logo}
\caption{武汉理工大学校徽（浮动图片）}
\label{fig:logo}
\end{figure}
\footnotemark[1]
\subsection{如何张贴源码？}
采用listing宏包可以完成代码的张贴，在控制文件导演区可以更改listing的设置
来符合MATLAB,Python，C++等不同语言的需求。
\begin{lstlisting}[language=C]
int main(int argc, char ** argv)
{
printf("Hello world!\n");
return 0;
}
\end{lstlisting}
\section{进阶功能}
\subsection{文献管理}
文献管理使用Bib\TeX ，可以从Google Scholar导出外文图书期刊等信息，从NoteExpress导出中文图书和期刊\cite{刘运来-199}。导出的信息基本格式类似于：
\begin{verbatim}
@article{
马晓丽-200,
   Author = {马晓丽},
   Title = {字体艺术的现代传承——有感于《字体设计》课程教育},
   Journal = {湖北成人教育学院学报},
   Volume = {19},
   Number = {2},
   Pages = {184-186},
   Abstract = {字体作为视觉传达设计中的一种符号文化，起着人与文化交流沟通的作用，是平面视觉传达设计的重要手段，这一点与图形的作用相通。汉字是代表中国文化的符号文字。因此，我们应该认真的研究它，从而发掘更多的造型方法，更深入地利用汉字来进行平面视觉传达设计。作为中国文化的继承者，我们应该自觉把文字艺术传承下去，创作出更多更好的设计造型。},
   Keywords = {字体; 视觉传达; 现代传承; 民族文化},
   Year = {2013} }
\end{verbatim}
如果需要引用该文献，可以直接使用\verb|\cite{马晓丽-200}|的方法进行引用。文章最后会自动根据GBT7714-2015规范来列出这些文献。

关于文献中出现[出版地不详]等问题，根目录下的两个bst文件默认设置忽略出版地等信息。如不需要该设置，可将其删除，使用texlive自带的宏包进行控制。
\subsection{转为Word}
Microsoft Word is the last thing I want to use before I die.
--Knuth

将本文档转化为word文档可以先转化为图片，再将所有图片插入到word文档中。
\section{已知问题和未来发展}
\subsection{已知问题}
本模板未采用2016版规范的页边距设置，因为实在是办不到2.5CM顶部页边距加上2.6CM的页眉设置啊。目录格式尚未修改，正在学习。
\subsection{未来发展}
武汉理工大学本科生论文的未来发展还是需要各位用户的参与，如果每一个用户都能贡献出一点关于\LaTeX 模板的想法和意见，我相信几年之后武汉理工大学本科生论文模板会成为其他高校学习和借鉴的例子。同学当自强，让我们一起来丰富完善这个模板，如果你有很好的建议或者意见请发送到 thesis@tsaoyu.com
\subsection{官方认证}
到目前为止（\today ）没有武汉理工大学任何官方组织对于本模板的格式或者内容进行认证，这代表采用本模板进行的论文写作可能不被官方的论文系统接受。如在进行原创性（防抄袭）检测的时候，可能需要提供提供doc版本的论文。希望用户了解到这个潜在的风险，做好文件转换和备份的准备。本人不对任何由于使用本模板而导致的毕业论文纠纷承担任何责任！
\section{致谢}
感谢父母为我提供的良好的衣食条件，让我有精力投入到这项没有经济回报的项目中去。
感谢徐海祥老师为我定制的论文题目，这个题目让我有兴趣制作这个模板。感谢武汉理工大学博士与硕士论文作者Hu,Weiyi，我在本模板制作的过程中参考了前辈的思路的方法。我研究过的模板还包括：上海交通大学，清华大学，哈尔滨工业大学，以及中国科技大学。其中论文引用格式GBT7714-2005-BibTeX-Style是上海财经大学的Haixing Hu作品，本模板离不开这些有益的资源的支持。同样感谢正在使用这个模板的你，相信通过你们的使用和传播，这个模板会变得越来越完善。
%============= 参考文献 =====================
\addcontentsline{toc}{section}{参考文献}
\printbibliography
%\bibliography{bibfile}
\clearpage
\newpage
\appendix

%=============  致谢  ======================
\section*{致~~谢}
\addcontentsline{toc}{section}{致谢}
感谢父母为我提供的良好的衣食条件，让我有精力投入到这项没有经济回报的项目中去。
感谢徐海祥老师为我定制的论文题目，这个题目让我有兴趣制作这个模板。感谢武汉理工大学博士与硕士论文作者Hu,Weiyi，我在本模板制作的过程中参考了前辈的思路的方法。我研究过的模板还包括：上海交通大学，清华大学，哈尔滨工业大学，以及中国科技大学。其中论文引用格式GBT7714-2005-BibTeX-Style是上海财经大学的Haixing Hu作品，本模板离不开这些有益的资源的支持。同样感谢正在使用这个模板的你，相信通过你们的使用和传播，这个模板会变得越来越完善。
\end{document}
%%%%%%%%%% 结束 %%%%%%%%%%